\section{Cronología del caso Sullivan: lo que se anticipaba en clase y el estado procesal final}

La clase utiliza el incidente de Uber de 2016 para discutir el bug bounty, la divulgación y la regulation by enforcement. Para mantener la precisión didáctica, es necesario leerlo en cuatro niveles: el relato personal de Sullivan en clase; la teoría del gobierno en la acusación y el juicio; las resoluciones que ya dictaron el jurado y los tribunales; y los controles de ingeniería que este texto extrae de todo ello. Ningún nivel puede sustituir a otro.

\subsection{Hitos procesales de 2016--2026}

El DOJ anunció en 2022 que el jurado emitió un veredicto de culpabilidad por obstrucción de un procedimiento de la FTC y por misprision of felony; en 2023 el tribunal impuso tres años de probation. Cuando se grabó la clase, Sullivan discutió el \term{nexus}: el requisito de vinculación que debe existir, en el contexto de la obstruction, entre la conducta imputada y el procedimiento oficial, y anticipó que el fallo Fischer de la Suprema Corte de 2024 podría favorecer la apelación. Posteriormente, el Noveno Circuito confirmó la condena y publicó una opinión modificada; la Suprema Corte denegó el \term{certiorari} (solicitud de revisión del expediente por el tribunal superior), es decir, se negó a conocer de esa solicitud de apelación, lo cual no significa que la Suprema Corte haya respaldado por separado cada cuestión de fondo.

\lecturefigure{10-current-legal-timeline.png}{La pending appeal mencionada en clase quedó sustituida por los resultados procesales de los tribunales en 2025--2026.}{DOJ; Ninth Circuit No. 23-927; Supreme Court docket 25-1082; redibujo conceptual.}

\begin{knowledgebox}{Lectura de la figura: separar hechos, opiniones y resoluciones}
2016 es el incidente de fondo; la jury conviction de 2022 y la sentencing de 2023 son resultados del juicio; la Ninth Circuit opinion de 2025 es la resolución de apelación; la cert denial de 2026 cierra esta ronda de solicitud de revisión ante la Suprema Corte. Lo que se afirma en clase sobre el nexus, el harmless error o el impacto en la industria corresponde a la postura procesal de enero de 2025. Las notas actuales pueden explicar sus implicaciones de gobernanza, pero no pueden seguir diciendo «apelación pendiente» ni presentar las predicciones del ponente como conclusiones del tribunal.
\end{knowledgebox}

\begin{warningbox}{El caso no es una prueba jurídica general para los VDP}
El caso involucra a una empresa concreta, una investigación de la FTC, los hechos del incidente, comunicaciones y registros. De un solo caso penal no puede deducirse que «pagar un bounty siempre es ilegal» ni que «firmar un NDA siempre es encubrimiento», como tampoco que «basta con tener un VDP para no incurrir en responsabilidad». Lo reutilizable en ingeniería son los mecanismos de clasificación, escalamiento, evidencia, divulgación y toma de decisiones conjunta; las conclusiones jurídicas concretas deben remitirse a la jurisdicción y a los hechos.
\end{warningbox}

\teachervoice{Sullivan veía en clase a Fischer y al nexus como una oportunidad importante para la apelación. La forma correcta de conservar la teacher voice es registrar «por qué lo juzgaba así en ese momento» y, al mismo tiempo, dejar claro que los resultados judiciales posteriores sustituyeron el estado pending; la fidelidad no consiste en congelar un momento histórico como si fuera un hecho actual.}

\begin{importantbox}{Cuatro etiquetas para las afirmaciones}
Al redactar un case study, distinga de forma explícita: \textbf{speaker account} (lo que afirma el ponente), \textbf{government allegation} (lo que sostiene el gobierno), \textbf{adjudicated fact/outcome} (lo que resolvieron el jurado o los tribunales) y \textbf{engineering judgment} (lo que recomienda este texto). Esto reduce las conjeturas sobre las motivaciones de las personas y permite al lector saber qué parte puede adoptarse directamente en el diseño de sistemas.
\end{importantbox}

\subsection{Resumen de la sección}

El valor didáctico de esta cronología está en corregir la versión anterior: al 11 de agosto de 2026, la apelación ya no está pending. Más importante aún, una organización no puede tratar el desenlace jurídico como una variable externa que solo importa después del incidente; el registro de los hechos, la asignación de responsabilidades y las vías de divulgación deben diseñarse antes de que ocurra el incidente.
